% Auto-generated by review/analyses/tissue_specific_features.py
% Top-20 features per tissue (Reviewer R3.4)
\begingroup
\scriptsize
\setlength{\tabcolsep}{4pt}
\renewcommand{\arraystretch}{1.12}
\begin{longtable}{@{}l r l p{0.55\linewidth} l r@{}}
\caption{\textbf{Top 20 developmental-stage features per tissue.} Genes ranked by aggregate LightGBM importance across the K$-$1 binary classifiers within the calibrated ordinal model for each tissue. Functional descriptions are drawn from mygene.info (\url{https://mygene.info/}) and UniProt (\url{https://rest.uniprot.org/}); placeholder entries equal the gene symbol indicate that no curated description was returned by either source. ``Peak stage'' is the developmental stage with the highest mean $\log_{2}$(TPM$+$1); ``FC'' is the peak minus trough in the same units. \label{tab:S7}} \\
\toprule
Tissue & Rank & Gene & Function & Peak stage & FC \\
\midrule
\endfirsthead
\multicolumn{6}{c}{\textit{Table~\ref{tab:S7} continued from previous page}} \\
\toprule
Tissue & Rank & Gene & Function & Peak stage & FC \\
\midrule
\endhead
\midrule
\multicolumn{6}{r}{\textit{continued on next page}} \\
\endfoot
\bottomrule
\endlastfoot
Muscle & 1 & \textit{SORCS2} & VPS10 domain-containing receptor SorCS2. & Infant & 2.87 \\
 & 2 & \textit{ARHGAP36} & GTPase activator for the Rho-type GTPases by converting them to an inactive GDP-bound state. & Infant & 4.76 \\
 & 3 & \textit{FGFRL1} & Has a negative effect on cell proliferation. & Infant & 3.98 \\
 & 4 & \textit{IGSF1} & Immunoglobulin superfamily member 1. & Infant & 4.04 \\
 & 5 & \textit{LOC100519808} & TRPM8 channel associated factor 1. & Infant & 3.90 \\
 & 6 & \textit{ACTC1} & Actin alpha cardiac muscle 1. & Infant & 10.14 \\
 & 7 & \textit{LOC102158612} & LOC102158612 & Post-pubertal & 4.89 \\
 & 8 & \textit{PPP1R14C} & Inhibitor of PPP1CA. & Early childhood & 4.61 \\
 & 9 & \textit{ACHE} & Carboxylic ester hydrolase. & Infant & 2.67 \\
 & 10 & \textit{LOC110257362} & LOC110257362 & Infant & 3.17 \\
 & 11 & \textit{TRDN} & Contributes to the regulation of lumenal Ca2+ release via the sarcoplasmic reticulum calcium release channels RYR1 and RYR2, a key step in triggering skeletal and heart muscle contraction. Required for normal organization of the triad junction, where T-tubules and the. & Post-pubertal & 2.09 \\
 & 12 & \textit{DLK1} & Delta like non-canonical Notch ligand 1. & Infant & 6.34 \\
 & 13 & \textit{ENSSSCG00000042876} & ENSSSCG00000042876 & Post-pubertal & 1.43 \\
 & 14 & \textit{SRPX} & Sushi repeat containing protein X-linked. & Infant & 4.53 \\
 & 15 & \textit{WDR1} & WD repeat-containing protein 1. & Infant & 2.18 \\
 & 16 & \textit{NRK} & Nik related kinase. & Infant & 5.38 \\
 & 17 & \textit{SLC22A16} & Facilitative organic cation transporter that mediates the transport of carnitine as well as the polyamine spermidine. Mediates the partially Na(+)-dependent bidirectional transport of carnitine. & Post-pubertal & 4.18 \\
 & 18 & \textit{ENSSSCG00000044729} & ENSSSCG00000044729 & Post-pubertal & 1.25 \\
 & 19 & \textit{ETV4} & ETS variant transcription factor 4. & Infant & 2.93 \\
 & 20 & \textit{GKAP1} & G kinase anchoring protein 1. & Adult & 1.38 \\
\midrule
Brain & 1 & \textit{ENSSSCG00000046117} & ENSSSCG00000046117 & Post-pubertal & 0.60 \\
 & 2 & \textit{CCDC9B} & Coiled-coil domain containing 9B. & Post-pubertal & 1.93 \\
 & 3 & \textit{SAMD5} & Sterile alpha motif domain-containing protein 5. & Post-pubertal & 1.09 \\
 & 4 & \textit{ENSSSCG00000038527} & Uncharacterized protein. & Adult & 0.72 \\
 & 5 & \textit{ENSSSCG00000045405} & ENSSSCG00000045405 & Post-pubertal & 0.53 \\
 & 6 & \textit{HMGCS2} & Catalyzes the condensation of acetyl-CoA with acetoacetyl-CoA to form HMG-CoA. & Infant & 0.22 \\
 & 7 & \textit{LATS1} & non-specific serine/threonine protein kinase. & Post-pubertal & 1.51 \\
 & 8 & \textit{GLIS3} & GLIS family zinc finger 3. & Post-pubertal & 1.33 \\
 & 9 & \textit{LOC100158196} & Electron carrier protein. The oxidized form of the cytochrome c heme group can accept an electron from the heme group of the cytochrome c1 subunit of cytochrome reductase. & Infant & 1.68 \\
 & 10 & \textit{NOTUM} & Notum, palmitoleoyl-protein carboxylesterase. & Pre-pubertal & 1.50 \\
 & 11 & \textit{SLC6A20} & Transporter. & Post-pubertal & 3.44 \\
 & 12 & \textit{FIBCD1} & Fibrinogen C domain containing 1. & Infant & 4.02 \\
 & 13 & \textit{LAMA1} & Anti-angiogenic and anti-tumor peptide that inhibits endothelial cell migration, collagen-induced endothelial tube morphogenesis and blood vessel growth in the chorioallantoic membrane. Blocks endothelial cell adhesion to fibronectin and type I collagen. & Post-pubertal & 1.34 \\
 & 14 & \textit{LRRC71} & Leucine rich repeat containing 71. & Adult & 1.04 \\
 & 15 & \textit{SASH1} & Is a positive regulator of NF-kappa-B signaling downstream of TLR4 activation. It acts as a scaffold molecule to assemble a molecular complex that includes TRAF6, MAP3K7, CHUK and IKBKB, thereby facilitating NF-kappa-B signaling activation. & Post-pubertal & 2.54 \\
 & 16 & \textit{LOC100516957} & Adhesion G protein-coupled receptor E2. & Pre-pubertal & 0.13 \\
 & 17 & \textit{CCDC170} & Coiled-coil domain containing 170. & Adult & 2.76 \\
 & 18 & \textit{CEP162} & Centrosomal protein of 162 kDa. & Post-pubertal & 1.49 \\
 & 19 & \textit{ACAP3} & GTPase-activating protein for the ADP ribosylation factor family. & Infant & 0.53 \\
 & 20 & \textit{ENSSSCG00000033041} & ENSSSCG00000033041 & Adult & 2.43 \\
\midrule
Liver & 1 & \textit{SATB2} & DNA-binding protein SATB. & Post-pubertal & 0.47 \\
 & 2 & \textit{CDH7} & Cadherins are calcium-dependent cell adhesion proteins. & Post-pubertal & 0.99 \\
 & 3 & \textit{ALDH18A1} & Bifunctional enzyme that converts glutamate to glutamate 5-semialdehyde, an intermediate in the biosynthesis of proline, ornithine and arginine. & Infant & 3.42 \\
 & 4 & \textit{ENSSSCG00000036118} & ENSSSCG00000036118 & Early childhood & 2.12 \\
 & 5 & \textit{ENSSSCG00000050649} & ENSSSCG00000050649 & Infant & 5.87 \\
 & 6 & \textit{COL1A1} & Collagen type I alpha 1 chain. & Infant & 1.90 \\
 & 7 & \textit{EIF3G} & RNA-binding component of the eukaryotic translation initiation factor 3 (eIF-3) complex, which is required for several steps in the initiation of protein synthesis. The eIF-3 complex associates with the 40S ribosome and facilitates the recruitment of eIF-1, eIF-1A,. & Early childhood & 0.96 \\
 & 8 & \textit{CD320} & CD320 molecule. & Infant & 2.03 \\
 & 9 & \textit{ACSF2} & Acyl-CoA synthases catalyze the initial reaction in fatty acid metabolism, by forming a thioester with CoA. Has some preference toward medium-chain substrates. & Infant & 1.92 \\
 & 10 & \textit{LOC100519366} & Cytochrome c oxidase subunit 6B1. & Post-pubertal & 1.79 \\
 & 11 & \textit{FAM169B} & Protein FAM169B. & Early childhood & 2.63 \\
 & 12 & \textit{DEK} & Involved in chromatin organization. & Adult & 1.10 \\
 & 13 & \textit{ENSSSCG00000033368} & ENSSSCG00000033368 & Infant & 1.28 \\
 & 14 & \textit{TK1} & Thymidine kinase. & Infant & 2.59 \\
 & 15 & \textit{SYN3} & Synapsin III. & Post-pubertal & 0.58 \\
 & 16 & \textit{RPL35A} & Component of the large ribosomal subunit. The ribosome is a large ribonucleoprotein complex responsible for the synthesis of proteins in the cell. & Early childhood & 1.42 \\
 & 17 & \textit{PTCHD1} & Required for the development and function of the thalamic reticular nucleus (TRN), a part of the thalamus that is critical for thalamocortical transmission, generation of sleep rhythms, sensorimotor processing and attention. Can bind cholesterol in vitro. & Post-pubertal & 0.44 \\
 & 18 & \textit{RARRES1} & Retinoic acid receptor responder 1. & Post-pubertal & 2.11 \\
 & 19 & \textit{GARNL3} & GTPase-activating Rap/Ran-GAP domain-like protein 3. & Post-pubertal & 2.44 \\
 & 20 & \textit{RAI2} & Retinoic acid induced 2. & Early childhood & 1.20 \\
\midrule
Lung & 1 & \textit{CCN5} & May play an important role in modulating bone turnover. Promotes the adhesion of osteoblast cells and inhibits the binding of fibrinogen to integrin receptors. & Post-pubertal & 2.71 \\
 & 2 & \textit{L1CAM} & Neural cell adhesion molecule involved in the dynamics of cell adhesion and in the generation of transmembrane signals at tyrosine kinase receptors. During brain development, critical in multiple processes, including neuronal migration, axonal growth and fasciculation, and. & Post-pubertal & 1.74 \\
 & 3 & \textit{GZMM} & Cleaves peptide substrates after methionine, leucine, and norleucine. Physiological substrates include EZR, alpha-tubulins and the apoptosis inhibitor BIRC5/Survivin. & Post-pubertal & 2.44 \\
 & 4 & \textit{HMGN1} & High mobility group nucleosome binding domain 1. & Infant & 0.98 \\
 & 5 & \textit{IGF2BP2} & RNA-binding factor that recruits target transcripts to cytoplasmic protein-RNA complexes (mRNPs). This transcript 'caging' into mRNPs allows mRNA transport and transient storage. & Infant & 1.57 \\
 & 6 & \textit{DSC2} & Desmocollin 2. & Infant & 1.21 \\
 & 7 & \textit{LOC102158691} & Immunoglobulin V-set domain-containing protein. & Post-pubertal & 4.04 \\
 & 8 & \textit{ADAM22} & ADAM metallopeptidase domain 22. & Post-pubertal & 0.62 \\
 & 9 & \textit{IGSF1} & Immunoglobulin superfamily member 1. & Infant & 0.80 \\
 & 10 & \textit{ENSSSCG00000012668} & Putative odorant or sperm cell receptor. & Post-pubertal & 0.02 \\
 & 11 & \textit{ENSSSCG00000041871} & ENSSSCG00000041871 & Post-pubertal & 0.01 \\
 & 12 & \textit{LOC100125545} & LOC100125545 & Early childhood & 0.36 \\
 & 13 & \textit{LOC110257692} & Putative odorant or sperm cell receptor. & Infant & 0.06 \\
 & 14 & \textit{ARHGAP36} & GTPase activator for the Rho-type GTPases by converting them to an inactive GDP-bound state. & Infant & 0.18 \\
 & 15 & \textit{LOC100511059} & Putative odorant or sperm cell receptor. & Early childhood & 0.03 \\
 & 16 & \textit{LOC100513516} & Putative odorant or sperm cell receptor. & Early childhood & 0.03 \\
 & 17 & \textit{OR13H1} & Putative odorant or sperm cell receptor. & Early childhood & 0.04 \\
 & 18 & \textit{LOC100514087} & Putative odorant or sperm cell receptor. & Early childhood & 0.03 \\
 & 19 & \textit{ENSSSCG00000048571} & ENSSSCG00000048571 & Early childhood & 0.05 \\
 & 20 & \textit{ENSSSCG00000035001} & ENSSSCG00000035001 & Pre-pubertal & 0.02 \\
\midrule
Blood & 1 & \textit{RPLP1} & Plays an important role in the elongation step of protein synthesis. & Infant & 3.14 \\
 & 2 & \textit{ENSSSCG00000033734} & ENSSSCG00000033734 & Early childhood & 8.93 \\
 & 3 & \textit{ENSSSCG00000042133} & ENSSSCG00000042133 & Adult & 3.05 \\
 & 4 & \textit{USF2} & Transcription factor that binds to a symmetrical DNA sequence (E-boxes) (5'-CACGTG-3') that is found in a variety of viral and cellular promoters. & Infant & 3.10 \\
 & 5 & \textit{BOLA2B} & Acts as a cytosolic iron-sulfur (Fe-S) cluster assembly factor that facilitates [2Fe-2S] cluster insertion into a subset of cytosolic proteins. Acts together with the monothiol glutaredoxin GLRX3. & Infant & 3.35 \\
 & 6 & \textit{TMEM158} & Transmembrane protein 158. & Post-pubertal & 1.95 \\
 & 7 & \textit{CCR6} & C-C motif chemokine receptor 6. & Post-pubertal & 2.33 \\
 & 8 & \textit{BPIFB4} & BPI fold containing family B member 4. & Post-pubertal & 0.86 \\
 & 9 & \textit{ENSSSCG00000045920} & ENSSSCG00000045920 & Infant & 3.56 \\
 & 10 & \textit{PKP2} & Plakophilin-2. & Adult & 1.39 \\
 & 11 & \textit{XPO7} & Exportin 7. & Infant & 4.71 \\
 & 12 & \textit{ENSSSCG00000051365} & ENSSSCG00000051365 & Adult & 3.76 \\
 & 13 & \textit{ENSSSCG00000045536} & ENSSSCG00000045536 & Adult & 4.10 \\
 & 14 & \textit{TLL1} & Metalloendopeptidase. & Adult & 0.27 \\
 & 15 & \textit{NSMCE4A} & Component of the SMC5-SMC6 complex, that promotes sister chromatid alignment after DNA damage and facilitates double-stranded DNA breaks (DSBs) repair via homologous recombination between sister chromatids. & Post-pubertal & 1.82 \\
 & 16 & \textit{NT5M} & 5',3'-nucleotidase, mitochondrial. & Post-pubertal & 1.21 \\
 & 17 & \textit{MIR9796} & microRNAs (miRNAs) are short (20-24 nt) non-coding RNAs that are involved in post-transcriptional regulation of gene expression in multicellular organisms by affecting both the stability and translation of mRNAs. miRNAs are transcribed by RNA polymerase II as part of capped and. & Post-pubertal & 0.04 \\
 & 18 & \textit{ENSSSCG00000031003} & Uncharacterized protein. & Post-pubertal & 4.29 \\
 & 19 & \textit{LOC100626120} & Ig-like domain-containing protein. & Post-pubertal & 1.33 \\
 & 20 & \textit{ENSSSCG00000034023} & ENSSSCG00000034023 & Early childhood & 0.55 \\
\end{longtable}
\endgroup
